\documentclass[aspectratio=169]{beamer}
\usepackage{ctex,hyperref}
\usepackage[T1]{fontenc}
\usepackage{latexsym,amsmath,xcolor,multicol,booktabs,array}
\usepackage{graphicx}
\usepackage{SUFE}

\author{汇报人：吴优\quad 学号：2025213385}
\title{基于去随机化Model-X Knockoff的电商GMV多粒度受控维度选择研究}
\subtitle{统一基准、组件消融与多粒度受控选择}
\institute{上海财经大学·统计与数据科学学院\\指导教师：张吕欧}
\date{开题报告·2026年9月}

\AtBeginSection[]{}
\AtBeginSubsection[]{}
\setbeamertemplate{navigation symbols}{}
\setbeamerfont{frame number}{parent=author in head/foot}
\graphicspath{{figures/}{pic/}}
\definecolor{deepred}{RGB}{124,67,76}
\definecolor{deepgreen}{RGB}{40,114,113}

\begin{document}

\kaishu
\begin{frame}
  \titlepage
  \begin{figure}
    \centering
    
\includegraphics[width=0.25\linewidth]{sufe.png}
  \end{figure}
\end{frame}

\begin{frame}{汇报提纲}
  \tableofcontents
\end{frame}

\section{研究背景}

\begin{frame}{研究背景：经营看板的维度过载}
  \begin{block}{现实问题}
    GMV看板同时包含交易规模、价格成本、供给、支付、履约、口碑、地域和营销等大量指标。指标越多，越难判断哪些变量真正提供了独立信息。
  \end{block}
  \begin{columns}[T]
    \begin{column}{0.49\textwidth}
      \textbf{常用做法}
      \begin{itemize}
        \item 业务经验筛选
        \item 边际相关与逐项检验
        \item 树模型重要性排序
      \end{itemize}
    \end{column}
    \begin{column}{0.49\textwidth}
      \textbf{主要不足}
      \begin{itemize}
        \item 高度相关代理重复入选
        \item 多重比较累积误报
        \item 重要性排序缺少错误比例控制
      \end{itemize}
    \end{column}
  \end{columns}
  \begin{alertblock}{研究目标}
    在31项候选维度中筛选对下一期GMV具有增量预测信息的少数指标，并形成可调数量、可复核的监控集合。
  \end{alertblock}
\end{frame}

\begin{frame}{研究问题与预期贡献}
  \begin{enumerate}
    \item 如何从多表交易数据构造不重复计算GMV的卖家面板？
    \item 如何在混合分布、高度相关变量中控制目标错误发现率？
    \item 如何在确认集合内通过参数$K$控制最终展示数量？
    \item 月度结论在周度、卖家内变化、业务组和非线性模型下是否稳定？
    \item 与标准筛选方法相比，主流程的增量价值及各组件作用是什么？
  \end{enumerate}
  \vspace{0.3cm}
  \begin{block}{预期贡献}
    构建“受控确认—公平比较—参数化展示—多粒度复核”的统一框架，严格区分预测精度、选择稳定性、错误率控制与因果结论。
  \end{block}
\end{frame}

\begin{frame}{相关研究与方法依据}
  \begin{columns}[T]
    \begin{column}{0.49\textwidth}
      \textbf{受控变量选择}
      \begin{itemize}
        \item BH错误发现率控制 [1]
        \item Knockoff过滤器 [2]
        \item Model-X Knockoff [3]
        \item 去随机化Knockoff [4]
        \item e-BH与e-value [5]
      \end{itemize}
    \end{column}
    \begin{column}{0.49\textwidth}
      \textbf{生成、预测与解释}
      \begin{itemize}
        \item MVR最小可重构性 [6]
        \item Deep Knockoffs [7]
        \item 深度核两样本检验 [8]
        \item XGBoost [9]
        \item TreeSHAP [10]
      \end{itemize}
    \end{column}
  \end{columns}
  \begin{alertblock}{研究空白}
    现有方法较少把受控发现、统一预算基准、组件消融、展示数量和月周粒度复核纳入同一经营指标治理流程。
  \end{alertblock}
\end{frame}

\section{数据与设计}

\begin{frame}{Olist数据与样本构造}
  \begin{columns}[T]
    \begin{column}{0.54\textwidth}
      \textbf{实际读取10张CSV}
      \begin{itemize}
        \item 订单与商品明细
        \item 支付与评价
        \item 商品、卖家与顾客
        \item 地理位置
        \item 营销线索与成交线索
      \end{itemize}
    \end{column}
    \begin{column}{0.44\textwidth}
      \textbf{数据质量原则}
      \begin{itemize}
        \item 自然粒度先聚合后连接
        \item GMV与明细逐笔对账
        \item 独立顾客使用唯一顾客标识
        \item 排除一项零方差变量
      \end{itemize}
    \end{column}
  \end{columns}
  \begin{table}
    \centering
    \small
    \begin{tabular}{lrrrr}
      \toprule
      口径 & 面板行 & 卖家 & 期间数 & 下一期零值率 \\
      \midrule
      月度主分析 & 14,677 & 2,793 & 19个月 & 25.86\% \\
      周度稳健性 & 33,190 & 2,834 & 83周 & 40.46\% \\
      \bottomrule
    \end{tabular}
  \end{table}
\end{frame}

\begin{frame}{31项候选维度与前瞻目标}
  \begin{columns}[T]
    \begin{column}{0.49\textwidth}
      \textbf{候选信息}
      \begin{itemize}
        \item 价格成本：价格、运费
        \item 商品结构：尺寸、重量、图片、文本
        \item 支付口碑：分期、支付方式、评分
        \item 履约：审核、承运、送达、承诺偏差
        \item 规模覆盖：订单、件数、顾客、品类、地域
        \item 营销：来源频率、成交线索匹配
      \end{itemize}
    \end{column}
    \begin{column}{0.49\textwidth}
      \textbf{预测目标}
      \begin{equation*}
        Y_{s,t+1}=\log(1+\mathrm{GMV}_{s,t+1})
      \end{equation*}
      \begin{itemize}
        \item 月度：本月预测次月
        \item 周度：本周预测下一周
        \item 下一期无成交时目标记为0
        \item 避免同周期GMV定义泄漏
      \end{itemize}
    \end{column}
  \end{columns}
\end{frame}

\section{研究方法}

\begin{frame}{统一方法框架}
  \begin{center}
    \large
    原始多表数据
    $\rightarrow$ 自然粒度聚合
    $\rightarrow$ 卖家面板
    $\rightarrow$ Copula-MVR Knockoff
    $\rightarrow$ 60轮e-value
    $\rightarrow$ e-BH确认集合
    $\rightarrow$ 公平基准与参数$K$
  \end{center}
  \vspace{0.5cm}
  \begin{columns}[T]
    \begin{column}{0.49\textwidth}
      \textbf{确认性主线}
      \begin{itemize}
        \item 随机化秩高斯Copula
        \item Ledoit-Wolf协方差
        \item MVR二阶Knockoff
        \item Lasso反对称统计量
      \end{itemize}
    \end{column}
    \begin{column}{0.49\textwidth}
      \textbf{补充验证}
      \begin{itemize}
        \item 时间外机器学习预测
        \item TreeSHAP与固定效应
        \item 卖家间/内与业务组
        \item 深度交换性诊断
      \end{itemize}
    \end{column}
  \end{columns}
\end{frame}

\begin{frame}[shrink=8]{Copula-MVR：从混合边际到有效负对照}
  \begin{block}{秩高斯变换}
    \begin{equation*}
      Z_{ij}=\Phi^{-1}\left(\frac{r_{ij}-0.5}{n}\right)
    \end{equation*}
  \end{block}
  \begin{block}{二阶条件生成}
    \begin{equation*}
      \widetilde Z\mid Z\sim
      \mathcal N\left(
        Z-Z\widehat\Sigma^{-1}S,\;
        2S-S\widehat\Sigma^{-1}S
      \right)
    \end{equation*}
  \end{block}
  \begin{itemize}
    \item $\widehat\Sigma$：Ledoit-Wolf收缩协方差
    \item $S$：按MVR准则降低真变量与仿制变量可重构性
    \item 交换要求：$(X,\widetilde X)_{\mathrm{swap}(A)}\overset d=(X,\widetilde X)$
  \end{itemize}
  \begin{alertblock}{解释边界}
    $H_{0j}:Y\perp X_j\mid X_{-j}$；竞争统计量
    $W_j=|\widehat\beta_j|-|\widehat\beta_{j+p}|$。拒绝零假设只表示条件预测信息。
  \end{alertblock}
\end{frame}

\begin{frame}[shrink=10]{e-value聚合、e-BH与参数$K$}
  \begin{block}{第$m$轮变量$j$的e-value}
    \begin{equation*}
      e_j^{(m)}
      =p\frac{\mathbf 1\{W_j^{(m)}\ge T^{(m)}\}}
      {1+\sum_{k=1}^{p}\mathbf 1\{W_k^{(m)}\le-T^{(m)}\}}
    \end{equation*}
  \end{block}
  \begin{equation*}
    \bar e_j=\frac{1}{60}\sum_{m=1}^{60}e_j^{(m)},\qquad
    \bar e_{(k)}\ge\frac{p}{\alpha_{\mathrm{eBH}}k}
  \end{equation*}
  \begin{itemize}
    \item 单轮$\alpha_{\mathrm{kn}}=0.10$；最终$\alpha_{\mathrm{eBH}}=0.20$
    \item e-BH承担正式FDR结论，单轮频率只描述随机稳定性
    \item $K^*=\min\{K:C_K\ge90\%,\ \mathrm{RMSE}_K\le1.005\,\mathrm{RMSE}_{\min}\}$
  \end{itemize}
\end{frame}

\begin{frame}{稳健性与诊断体系}
  \begin{table}
    \centering
    \small
    \begin{tabular}{p{0.29\textwidth}p{0.29\textwidth}p{0.32\textwidth}}
      \toprule
      风险 & 方法 & 回答的问题 \\
      \midrule
      卖家内相关 & 均值—离差分层 & 信号来自跨卖家还是卖家自身变化？ \\
      共线代理 & 12业务组Knockoff & 哪类业务信息整体有效？ \\
      非高斯联合分布 & 深度生成+深度核MMD & 生成器是否满足联合交换性？ \\
      非线性重要性 & 轨迹聚合配对MLP & 非线性竞争能否恢复信号？ \\
      粒度依赖 & 周度完整重跑 & 月度结论是否依赖聚合尺度？ \\
      \bottomrule
    \end{tabular}
  \end{table}
\end{frame}

\begin{frame}{基准比较与组件消融设计}
  \begin{columns}[T]
    \begin{column}{0.48\textwidth}
      \textbf{公平比较协议}
      \begin{itemize}
        \item 选维数据截至2018年4月
        \item 统一变量预算：12项
        \item 统一预测器：同参数XGBoost
        \item 统一测试期：2018年5--7月
        \item 30组相同的70\%卖家子样本
        \item 比较RMSE、配对Jaccard与集合重合
      \end{itemize}
      \textbf{五种方法}
      \begin{itemize}
        \item 去随机化Knockoff、Elastic Net
        \item 稳定性选择、影子变量树、XGBoost-SHAP
      \end{itemize}
    \end{column}
    \begin{column}{0.48\textwidth}
      \textbf{组件消融}
      \begin{itemize}
        \item 边际：Copula / 原始高斯
        \item $S$矩阵：MVR / SDP / 等相关
        \item 统计量：Lasso / XGBoost
        \item 重复数：$M=1,5,10,20,40,60$
      \end{itemize}
      \textbf{已知真值模拟}
      \begin{itemize}
        \item 真实特征协方差，$n=4000,p=31$
        \item 10个真信号，20轮响应重复
        \item FDP、FDP超标率、Power、发现数
      \end{itemize}
    \end{column}
  \end{columns}
\end{frame}

\section{预分析结果}

\begin{frame}[shrink=8]{月度确认性结果}
  \begin{columns}[T]
    \begin{column}{0.48\textwidth}
      \includegraphics[width=\linewidth]{fig03_knockoff_evalues.png}
    \end{column}
    \begin{column}{0.50\textwidth}
      \textbf{$q=0.20$确认14项}
      \begin{itemize}
        \item 规模：订单数、件数、顾客数、品类数
        \item 价格：平均价格、平均运费
        \item 履约：承运、送达、预计偏差
        \item 内容与覆盖：标题、描述、高度、买家州、卖家州
      \end{itemize}
      \textbf{高稳定项：}10项单轮频率$\geq90\%$
    \end{column}
  \end{columns}
  \begin{alertblock}{阈值现象}
    $q=0.10$为0项、$q=0.20$为14项，说明低维Knockoff+存在离散门槛。
  \end{alertblock}
\end{frame}

\begin{frame}[shrink=10]{参数$K$与月度预测}
  \begin{columns}[T]
    \begin{column}{0.46\textwidth}
      \textbf{默认$K=10$}
      \begin{enumerate}
        \item 独立顾客数
        \item 成交订单数
        \item 平均商品价格
        \item 审核至承运天数
        \item 预计提前送达天数
        \item 买家州数量
        \item 下单至送达天数
        \item 成交件数
        \item 平均运费
        \item 商品标题长度
      \end{enumerate}
    \end{column}
    \begin{column}{0.52\textwidth}
      \begin{block}{时间外结果}
        \begin{itemize}
          \item 默认$K=10$：RMSE=2.385
          \item 全31维XGBoost：RMSE=2.388
          \item 相对变化：$-0.097\%$
          \item 维度压缩后性能基本保留
        \end{itemize}
      \end{block}
      \includegraphics[width=\linewidth]{fig10_predictive_models.png}
    \end{column}
  \end{columns}
\end{frame}

\begin{frame}[shrink=12]{统一基准、消融与模拟结果}
  \begin{columns}[T]
    \begin{column}{0.49\textwidth}
      \textbf{真实数据：相同12项预算}
      \begin{table}
        \scriptsize
        \begin{tabular}{lrr}
          \toprule
          方法 & RMSE & 稳定性\\
          \midrule
          去随机化Knockoff & 2.392 & 0.855\\
          Elastic Net & 2.390 & 0.873\\
          稳定性选择 & 2.390 & 0.814\\
          影子变量树 & 2.387 & 0.832\\
          XGBoost-SHAP & 2.383 & 0.686\\
          \bottomrule
        \end{tabular}
      \end{table}
      \begin{block}{解释}
        SHAP预测略优；Elastic Net与Knockoff稳定性区间重叠，Knockoff另承担FDR解释。
      \end{block}
    \end{column}
    \begin{column}{0.49\textwidth}
      \textbf{20轮模拟}
      \begin{table}
        \scriptsize
        \begin{tabular}{lrr}
          \toprule
          方法 & FDP & Power\\
          \midrule
          去随机化Knockoff & 0.054 & 1.000\\
          单次Knockoff & 0.145 & 1.000\\
          Elastic Net & 0.644 & 1.000\\
          稳定性选择 & 0.090 & 1.000\\
          影子变量树 & 0.566 & 1.000\\
          \bottomrule
        \end{tabular}
      \end{table}
      \begin{alertblock}{消融结论}
        重复数从1到60时集合仍有波动；原始高斯、替代$S$矩阵和XGBoost统计量均改变结果。
      \end{alertblock}
    \end{column}
  \end{columns}
\end{frame}

\begin{frame}[shrink=12]{周度稳健性结果}
  \begin{columns}[T]
    \begin{column}{0.48\textwidth}
      \includegraphics[width=\linewidth]{fig17_monthly_weekly_frequency.png}
    \end{column}
    \begin{column}{0.50\textwidth}
      \begin{itemize}
        \item 周度$q=0.20$入选18项
        \item 与月度共同入选11项
        \item Jaccard系数为0.524
        \item 周度高稳定项13项
        \item 周度18项XGBoost RMSE=2.421
        \item 全31维XGBoost RMSE=2.420
      \end{itemize}
    \end{column}
  \end{columns}
  \begin{block}{跨粒度核心}
    交易规模、平均价格与运费、承运准备、商品标题和区域覆盖在两种粒度下共同出现。
  \end{block}
\end{frame}

\begin{frame}{结构检验与负结果}
  \begin{table}
    \centering
    \small
    \begin{tabular}{lccp{0.36\textwidth}}
      \toprule
      分支 & 月度 & 周度 & 解释 \\
      \midrule
      卖家间层$q=0.20$ & 16项 & 12项 & 长期卖家差异包含主要信号 \\
      卖家内层$q=0.20$ & 0项 & 0项 & 短期动态证据不足 \\
      业务组$q=0.20$ & 0组 & 10组 & 周度组级信息更强 \\
      轨迹重要性$q=0.20$ & 0项 & 0项 & 主水平未恢复非线性发现 \\
      \bottomrule
    \end{tabular}
  \end{table}
  \begin{alertblock}{深度生成器准入}
    边际KS较小，但深度核MMD拒绝率100\%、分类器AUC为0.902/0.937；因此不作FDR声明。
  \end{alertblock}
\end{frame}

\section{计划与产出}

\begin{frame}[shrink=6]{研究创新点与方法主线}
  \begin{enumerate}
    \item \textbf{单一确认主线：}Copula-MVR、反对称$W$、e-value与e-BH共同承担正式发现。
    \item \textbf{确认—展示分离：}参数$K$只压缩确认池，不反向改变FDR阈值。
    \item \textbf{分层证据：}月度确认，周度、业务组、卖家间/内与深度诊断负责复核。
    \item \textbf{明确否决：}未通过交换性诊断的深度生成器不得用于FDR声明。
    \item \textbf{公平评价：}统一窗口、预算、预测器与测试期，并用消融和模拟解释组件作用。
  \end{enumerate}
\end{frame}

\begin{frame}{后续研究计划}
  \begin{table}
    \centering
    \small
    \begin{tabular}{lp{0.62\textwidth}}
      \toprule
      时间 & 工作内容 \\
      \midrule
      2026.09 & 完成研究设计、数据审计与预分析整理 \\
      2026.10 & 完善理论推导、模拟设计和稳健性讨论 \\
      2026.11 & 统一图表、核对文献与结果叙述 \\
      2026.12 & 完成论文初稿、导师反馈与结构修改 \\
      2027.01 & 完成格式终检、盲审稿与答辩材料 \\
      \bottomrule
    \end{tabular}
  \end{table}
\end{frame}

\section{参考文献}

\begin{frame}[shrink=22]{主要参考文献}
  \tiny
  \begin{columns}[T]
    \begin{column}{0.49\textwidth}
      \begin{enumerate}
        \item Benjamini \& Hochberg. Controlling the false discovery rate. JRSS B, 1995.
        \item Barber \& Candès. Controlling the false discovery rate via knockoffs. AoS, 2015.
        \item Candès et al. Panning for gold: Model-X knockoffs. JRSS B, 2018.
        \item Barber et al. Robust inference with knockoffs. AoS, 2020.
        \item Ren \& Barber. Derandomised knockoffs. JRSS B, 2024.
        \item Wang \& Ramdas. FDR control with e-values. JRSS B, 2022.
        \item Vovk \& Wang. E-values: calibration, combination and applications. AoS, 2021.
        \item Ledoit \& Wolf. A well-conditioned covariance estimator. JMA, 2004.
        \item Spector \& Janson. Powerful knockoffs via minimizing reconstructability. AoS, 2022.
        \item Bates et al. Metropolized knockoff sampling. JASA, 2021.
        \item Romano et al. Deep knockoffs. JASA, 2020.
        \item Dai \& Barber. Knockoff filter for group-sparse regression. ICML, 2016.
        \item Liu et al. Learning deep kernels for two-sample tests. ICML, 2020.
      \end{enumerate}
    \end{column}
    \begin{column}{0.49\textwidth}
      \begin{enumerate}
        \setcounter{enumi}{13}
        \item Tibshirani. Regression shrinkage and selection via the Lasso. JRSS B, 1996.
        \item Zou \& Hastie. Regularization and variable selection via Elastic Net. JRSS B, 2005.
        \item Breiman. Random forests. Machine Learning, 2001.
        \item Friedman. Gradient boosting machine. AoS, 2001.
        \item Chen \& Guestrin. XGBoost. KDD, 2016.
        \item Lundberg \& Lee. A unified approach to interpreting predictions. NeurIPS, 2017.
        \item Lundberg et al. Explainable AI for trees. Nature MI, 2020.
        \item Wooldridge. Econometric Analysis of Cross Section and Panel Data. MIT, 2010.
        \item Little \& Rubin. Statistical Analysis with Missing Data. Wiley, 2019.
        \item Hyndman \& Athanasopoulos. Forecasting: Principles and Practice, 2021.
        \item Olist. Brazilian E-Commerce Public Dataset. Kaggle, 2018.
        \item Olist. Marketing Funnel Dataset. Kaggle, 2018.
      \end{enumerate}
    \end{column}
  \end{columns}
\end{frame}

\end{document}
